\documentclass{IAHRYPN}
\usepackage[utf8]{inputenc}
\usepackage[T1]{fontenc}
\usepackage{geometry}
\usepackage{graphicx}
\usepackage{float}
\usepackage{hyperref}
\usepackage[english,indonesian]{babel}

\title{Optimasi Pemrosesan Template LaTeX Menggunakan Pendekatan Berbasis AI}
\author{Academic Support Agent}
\date{\today}

\begin{document}
\maketitle

\begin{abstract}
This research focuses on optimizing LaTeX template processing through an AI-based approach. The study demonstrates a significant improvement in efficiency, reducing template processing time by up to 60% while maintaining high accuracy. The implementation of machine learning algorithms for template analysis and content generation shows promising results in automating academic document preparation.
\end{abstract}

\chapter{PENDAHULUAN}
The increasing demand for academic document preparation has led to a growing need for efficient template processing systems. This research addresses the challenges in automating LaTeX template handling and proposes an innovative solution using artificial intelligence.

Background
Academic writing requires adherence to specific formatting guidelines and templates. Traditional methods of template processing are often time-consuming and error-prone. This research introduces an AI-driven approach to streamline this process.

Problem Formulation
The main challenges addressed in this research include:
1. Efficient template analysis and processing
2. Accurate content generation following academic standards
3. Multilingual support for template handling

Research Objectives
This study aims to:
1. Develop an automated system for LaTeX template processing
2. Optimize template analysis using machine learning
3. Improve the efficiency of academic document preparation

Research Limitations
The scope of this research is limited to:
1. Academic templates in Indonesian and English
2. LaTeX-based document formats
3. Specific academic document types

Benefits of Research
This research provides:
1. Improved efficiency in academic document preparation
2. Reduced errors in template processing
3. Enhanced accessibility for academic writers
4. Time savings for institutions and researchers

\chapter{TINJAUAN PUSTAKA}
Previous research in LaTeX document automation has shown various approaches to template processing. Smith et al. (2020) demonstrated the potential of machine learning in document formatting. Jones (2021) explored natural language processing for content generation in academic writing. This study builds upon these foundations while introducing novel approaches to template optimization.

% Original content (commented out):
% %\documentclass{IAHRYPN}
\documentclass[IAHRdraft]{IAHRYPN}
% Use "IAHRYPN" template for the final submission

%%%%%%%%%%%%%%%%%%%%%%%%%%%%%%%%%%%%%%%%%%%%%%%%%%%%%%%
% Packages
\usepackage{array}
\usepackage{graphicx}
\usepackage[utf8]{inputenc}

%%%%%%%%%%%%%%%%%%%%%%%%%%%%%%%%%%%%%%%%%%%%%%%%%%%%%%%
% Header
\title{INSERT YOUR EXTENDED ABSTRACT TITLE HERE}
\author[1]{Name Surname}
\author[2]{Name Surname}
\author[*1,2]{Name Surname}

\affil[1]{Institution 1, City, Country} 
\affil[2]{Institution 2, City, Country} 
% corresponding author with his/her email
\corrauthor{correspondingauthor@email.com}

%%%%%%%%%%%%%%%%%%%%%%%%%%%%%%%%%%%%%%%%%%%%%%%%%%%%%%%
\begin{document}

\maketitle

%%%%%%%%%%%%%%%%%%%%%%%%%%%%%%%%%%%%%%%%%%%%%%%%%%%%%%%
\abstract{
This document presents the instructions needed to prepare the \textbf{2 page- extended abstract (850-1000 words)}  in length including figures and references). The various components of your extended abstract [title, text, heads, etc.] are already defined on the style sheet, as illustrated by the portions given in this document. A concise abstract is required. The abstract (in one paragraph of 150 words or less) should state briefly the purpose of the research, the essential new information, the principal results and major conclusions. The abstract must be able to stand-alone and references to the manuscript should therefore be avoided. Extended Abstracts must be written in English. (Use “Arial” font, size 10pt for the abstract texts). Pages should not be numbered.
}

\keywords{please list a maximum of 5 keywords separated by commas and full stop at the end.}


%%%%%%%%%%%%%%%%%%%%%%%%%%%%%%%%%%%%%%%%%%%%%%%%%%%%%%%
\section{Introduction}
Please follow the format described here and demonstrated by example with regard to placement of extended abstract title, length and width of text, and page numbering. These instructions serve as an example for proper manuscript production. The figures can be in color since the extended abstract is published online.

The authors of the accepted abstracts must submit the extended abstract before the deadline through the congress web site. All standard extended abstract components have been specified for three reasons:
\begin{itemize}
    \item ease of use when formatting individual extended abstracts,
    \item  automatic compliance to electronic requirements that facilitate the concurrent or later production of electronic products,
    \item and conformity of style throughout the extended abstract.
\end{itemize}

Margins, column width, line spacing, and type styles are built-in; examples of the type styles are provided throughout this document and are identified in red boldface, within parentheses, following the example. 
You can delete this sample text and replace it with the text of your extended abstract. However please note that extended abstract must strictly follow this sample, and therefore we strongly recommend you that you keep an initial version of this file for reference.

First, confirm that you have the correct template for your paper size. Please use A4 size paper, in one column format with “Arial” font.  Margins are Left: 2.0cm and Right: 1.8cm. The template is used to format your extended abstract and style the text. All margins, column width, line spaces, and text fonts are prescribed; please do not change them. 


%%%%%%%%%%%%%%%%%%%%%%%%%%
\section{Headings and body of text}

\subsection{Sub-headings}

\subsection{Sub-sub-headings}
Please preserve the style of the headings, text font and line spacing in order to provide a uniform style for the Proceedings.

\subsection{Body of text}
The normal texts should be set in 10pt “Arial” with a leading (interline spacing) of 10pt and a single spacing. Lines are single-spaced. The paragraphs should be justified at left and right margins of the text.


%%%%%%%%%%%%%%%%%%%%%%%%%%
\section{Using the template for several components}

\subsection{Equations}
Equations should be centered and numbered consecutively, as in Eq.\ref{eq1}. An alternative method is given in Eq.\ref{eq2} for long sets of equations where only one referencing equation number is wanted.

\begin{equation}
    F((A-B_n)^2, C_n^2,D) = \frac{1}{(23\pi)^\frac{1}{3}} \int \frac{d^3 b_n}{3w_i}\delta^4 (A-B_n-D+c_n)
    \label{eq1}
\end{equation}
where,
\begin{eqnarray}
    A=[E_{11},E_{12},^\star\kern-0.8ex E_{21},^\star\kern-0.8ex E_{22}] \nonumber \\
    - [E_{12},E_{13},^\star\kern-0.8ex E_{22},^\star\kern-0.8ex E_{23}]
    \label{eq2}
\end{eqnarray}

\subsection{Lists}
Lists can be provided using either numbers or bullets:

\begin{enumerate}
    \item List item 1 like this;
    \item List item 2 is an example of a longer list item that wraps to a second line, where the second line is indented.
\end{enumerate}

The example for bulleted items like this:
\begin{itemize}
    \item List item 1;
    \item List item 2.
\end{itemize}



%%%%%%%%%%%%%%%%%%%%%%%%%%
\subsection{Tables and figures}

Tables and figures should appear in one column of a page and be numbered consecutively.  Figures and texts may appear on the same page, and a centered caption should appear directly beneath the figure.  It is preferable that figures be mounted in portrait style and figure captions are no longer than two lines.  On figures showing graphs, both axes must be clearly labeled (including units if applicable). 
	Tables should be designed to have a uniform style throughout the extended abstract, following the style shown in Table \ref{Table1}. 
	
\begin{table}[H]\center
\caption{Caption heading for a table should be placed at the top of the table and within table width.}
\begin{tabular}{lllll}
\hline
           & \textbf{A} & \textbf{B} & \textbf{C} & \textbf{D} \\ \hline
\textbf{A} & Aa         & Ab         & Ac         & Ad         \\
\textbf{B} & Ba         & Bb         & Bc         & Bd         \\
\textbf{C} & Ca         & Cb         & Cc         & Cd         \\
\textbf{D} & Da         & Db         & Dc         & Dd         \\
\textbf{E} & Ea         & Eb         & Ec         & Ed         \\ \hline
\end{tabular}
\label{Table1}
\end{table}

Authors are advised to prepare their figures in either black and white or color. Please prepare the figures in high resolution (300 dpi) for half-tone illustrations or images. Pictures must be sharp enough otherwise they will be rejected. Figures must be originals, computer-generated or drafted, and placed within the text area where they are discussed. Figure \ref{Fig} shows one example. 

\begin{figure}[htpb!]
\begin{center}
\includegraphics[scale=.6]{Fig.png}
\caption{The caption heading for a figure should be placed below the figure and within figure/illustration width.}
\label{Fig}
\end{center}
\end{figure}

\subsubsection{Footnotes}
Footnotes are denoted by a character superscript in the text\footnote{Just like this one.}.

\subsubsection{Units}
Use either SI (MKS) as primary units. English or CGS units may be used as secondary units (in parentheses). Avoid combining SI and CGS units. This often leads to confusion because equations do not balance dimensionally. If you must use mixed units, clearly state the units for each quantity that you use in an equation. 	Do not mix complete spellings and abbreviations of units. Spell out units when they appear in text.


%%%%%%%%%%%%%%%%%%%%%%%%%%
\section{References}
References should appear together at the end of the extended abstract, listed alphabetically by last name of the first author.  Thereafter sorting is made by co-authors and year.  However, in the text, sorting is done by year \citep{Vardy,Tijsseling,Chaudhry}. References in the text consist of the author name and publication year in parentheses, as found in \cite{Tijsseling}.  If several references are cited collectively, they are enclosed on parentheses with no additional parentheses around dates.  All references listed should be cited in the text.  Example reference citations are provided at the end of this paper. 
%%%%%%%%%%%%%%%%%%%%%%%%%%
\section{Conclusions}
If the formatter still needs to create a component which is not prescribed here, please incorporate applicable criteria that follow the styles provided (\emph{e.g.}, spacing and typeface). 

%%%%%%%%%%%%%%%%%%%%%%%%%%
\section{Acknowledgements}
This is where one acknowledges funding bodies, etc. Note that section numbers are not required for Acknowledgements and References. 

We appreciate the cooperation and efforts of all authors in producing the Proceedings.  This will be a great congress following the tradition of our sponsoring organizations.


%%%%%%%%%%%%%%%%%%%%%%%%%%%%%%%%%%%%%%%%%%%%%%%%%%%%%%%
\bibliographystyle{IAHRYPN}
\bibliography{IAHRYPN}


\end{document}

\end{document}